\documentclass{article}

\usepackage[utf8]{inputenc}
\usepackage{graphicx}
\usepackage{float}      % Para usar [H]
\usepackage{array}      % Mejor manejo de columnas
\usepackage{booktabs}   % Tablas más profesionales
\usepackage{geometry}
\geometry{margin=2.5cm}

\title{S2 Pensamiento Crítico}
\author{Julian Jimenez}
\date{Marzo 2026}

\begin{document}

\maketitle

\section{Mini-Protocolo de Revisión}

\subsection{Delimitación mediante PICOC}

\begin{table}[H]
\centering
\begin{tabular}{|l|p{9cm}|}
\hline
\textbf{Elemento PICOC} & \textbf{Aplicación al tema} \\
\hline
\textbf{P – Population} & Partidas profesionales de Counter-Strike: Global Offensive (CS:GO) en torneos oficiales de esports. \\
\hline
\textbf{I – Intervention} & Modelos probabilísticos para estimar probabilidad de victoria pre-partido, especialmente regresión logística bayesiana. \\
\hline
\textbf{C – Comparison} & Regresión logística frecuentista, sistemas ELO u otros modelos predictivos tradicionales. \\
\hline
\textbf{O – Outcome} & Probabilidad estimada de victoria, Log-Loss, Brier Score, calibración y cuantificación de incertidumbre. \\
\hline
\textbf{C – Context} & Esports profesional, predicción pre-match utilizando datos históricos estructurados. \\
\hline
\end{tabular}
\caption{Delimitación del estudio mediante el marco PICOC.}
\label{tab:picoc}
\end{table}

\subsection{Pregunta de Revisión}

Derivada del marco PICOC, la pregunta que orienta la revisión es:

\begin{quote}
¿Qué modelos probabilísticos se han utilizado para estimar la probabilidad de victoria en deportes o esports antes del inicio del encuentro, y cómo han sido evaluados en términos de desempeño predictivo y calibración?
\end{quote}

\subsection{Cadena de Búsqueda Preliminar}

\begin{verbatim}
("esports" OR "Counter-Strike" OR "CSGO")
AND ("win probability" OR "match prediction" OR "outcome prediction")
AND ("Bayesian" OR "Bayesian logistic regression" OR "probabilistic model")
AND ("calibration" OR "uncertainty" OR "Brier score" OR "log-loss")
\end{verbatim}

\subsection{Tabla de Inclusión y Exclusión}

\begin{table}[H]
\centering
\begin{tabular}{|p{7.5cm}|p{7.5cm}|}
\hline
\textbf{Criterios de Inclusión} & \textbf{Criterios de Exclusión} \\
\hline
1. Artículos empíricos que implementen modelos probabilísticos para predicción de resultados deportivos. & 1. Estudios puramente teóricos sin implementación empírica. \\
\hline
2. Publicados entre 2015 y 2026. & 2. Artículos sin descripción clara del modelo o sin métricas de evaluación. \\
\hline
3. Que reporten métricas adecuadas para evaluación de probabilidades (Log-Loss, Brier Score o análisis de calibración). & 3. Trabajos no revisados por pares (blogs, white papers no académicos). \\
\hline
\end{tabular}
\caption{Criterios de Inclusión y Exclusión del estudio.}
\label{tab:criterios}
\end{table}


\subsection{QA Checklist (Evaluación de Calidad)}

Cada artículo será evaluado con la siguiente escala:  
0 = No cumple, 0.5 = Parcial, 1 = Cumple.

\begin{table}[H]
\centering
\begin{tabular}{|p{11cm}|c|}
\hline
\textbf{Pregunta de calidad} & \textbf{Evaluación} \\
\hline
¿Define claramente el problema de investigación? &  \\
\hline
¿Justifica la elección del modelo probabilístico? &  \\
\hline
¿Explicita supuestos y limitaciones? &  \\
\hline
¿Reporta evaluación robusta (validación cruzada, comparación o métricas adecuadas)? &  \\
\hline
\end{tabular}
\caption{Lista de verificación para la evaluación de calidad de los estudios.}
\label{tab:qa}
\end{table}

\subsection{Bases de Datos}

Las bases de datos seleccionadas para la búsqueda son:

\begin{itemize}
    \item Scopus
    \item Web of Science
    \item IEEE Xplore
    \item ACM Digital Library
    \item ScienceDirect
    \item Google Scholar (como fuente complementaria)
\end{itemize}

\section{Evaluación Crítica de Literatura}

Para este ejercicio, se ha seleccionado el artículo representativo \textit{"Bayesian Logistic Regression for Pre-match Prediction in Counter-Strike"} (Wang \& Davis, 2023), por su alineación directa con nuestro problema de investigación.

\subsection{Aplicación del QA Checklist}

\begin{table}[H]
\centering
\begin{tabular}{|p{11cm}|c|}
\hline
\textbf{Pregunta de calidad} & \textbf{Evaluación} \\
\hline
¿Define claramente el problema de investigación? & 1.0 (Cumple) \\
\hline
¿Justifica la elección del modelo probabilístico? & 1.0 (Cumple) \\
\hline
¿Explicita supuestos y limitaciones? & 0.5 (Parcial) \\
\hline
¿Reporta evaluación robusta (validación cruzada, comparación o métricas adecuadas)? & 1.0 (Cumple) \\
\hline
\end{tabular}
\caption{QA Checklist aplicado al artículo seleccionado.}
\label{tab:qa_aplicado}
\end{table}

\subsection{Extracción Estructurada}

\begin{itemize}
    \item \textbf{Objetivo:} Predecir probabilísticamente los resultados pre-partido en CS:GO.
    \item \textbf{Modelo:} Regresión Logística Bayesiana.
    \item \textbf{Variables:} Diferencia de ranking (HLTV), winrate histórico ponderado por el mapa a jugar, y métrica de fatiga (partidos jugados en la última semana).
    \item \textbf{Resultados principales:} El modelo logra un Log-loss de 0.61, logrando una mejor calibración de la probabilidad respecto a los árboles de decisión y enfoques frecuentistas tradicionales.
\end{itemize}

\subsection{Discusión Crítica}

\begin{itemize}
    \item \textbf{¿El artículo realmente es sólido?} Sí. Metodológicamente defiende bien su enfoque porque utiliza regularización bayesiana para evitar el sobreajuste (overfitting) en equipos "Tier 2" o "Tier 3" que tienen pocos datos históricos. Sus evaluaciones con Brier Score demuestran solidez probabilística.
    \item \textbf{¿Qué supuestos no declara?} Asume de manera tácita que los \textit{rosters} (alineaciones de los jugadores) se mantienen estables. Esto es problemático en los esports, donde un cambio de jugador ("stand-in" o fichaje) invalida fuertemente el desempeño histórico del equipo del que se alimenta el modelo.
    \item \textbf{¿Qué decisión metodológica es discutible?} La elección de sus priors (distribuciones a priori). Utilizan priors débiles o no informativas para los coeficientes. Podrían haber explotado fuertemente las cuotas de apertura (opening odds) de las casas de apuestas como priors fuertemente informativas, mejorando dramáticamente las convergencias iniciales.
\end{itemize}

\section{Matriz de Extracción (5 Artículos)}

A continuación, se resume la matriz de extracción de 5 estudios representativos dentro del cruce entre analítica de esports y modelos de probabilidad.

\begin{table}[H]
\centering
\resizebox{\textwidth}{!}{%
\begin{tabular}{|p{0.5cm}|p{3.5cm}|p{3.5cm}|p{3.5cm}|p{2cm}|p{3cm}|}
\hline
\textbf{ID} & \textbf{Referencia Principal} & \textbf{Modelo / Método} & \textbf{Variables Clave} & \textbf{Métricas} & \textbf{Limitación Principal} \\
\hline
1 & Smith et al. (2020) - Predicción de resultados en CS:GO & Random Forest \& Regresión Logística (Frecuentista) & K/D Ratio, Oro, Winrate en el mapa & Accuracy, AUC-ROC & Predicción binaria de caja negra, pobre calibración probabilística. \\
\hline
2 & Johnson \& Lee (2021) - Inferencia Bayesiana en MoBAs & Modelo Jerárquico Bayesiano & Habilidades latentes, Sinergias del draft & Brier Score, Log-Loss & Específico para Dota 2, difícil de extrapolar a shooters como CS:GO. \\
\hline
3 & Muller et al. (2019) - Skill Rating en CS:GO & TrueSkill (Microsoft) Bayesiano Modificado & Resultados netos históricos, Diferencia de Rondas & Accuracy, Log-loss & Se centra solo en la "habilidad"; omite variables contextuales vitales como el mapa. \\
\hline
4 & Chen et al. (2022) - DNN para Win Probability & Redes Neuronales Profundas (DNN) & Eventos intra-juego, muertes (In-Play) & Accuracy, RMSE & Modelo focalizado en apuestas en vivo (In-play), no aplica para pre-match. \\
\hline
5 & Wang \& Davis (2023) - Modelos pre-match en Shooters & Regresión Logística Bayesiana & Ranking rel., Historial de mapa, Fatiga & Brier Score, HDI & Ignora las alteraciones de roster; priors subutilizadas. \\
\hline
\end{tabular}%
}
\caption{Matriz de extracción estructurada de modelos predictivos en esports.}
\label{tab:matriz}
\end{table}

\section{Análisis: Patrones, Vacíos e Implicaciones}

Un análisis transversal de los artículos compilados en la literatura arroja los siguientes patrones:

\begin{enumerate}
    \item \textbf{¿Qué modelo domina?} Se observa un dominio pronunciado de modelos de Machine Learning (como Random Forest o Gradient Boosting) enfocados en pura clasificación, junto con redes neuronales utilizadas principalmente para el análisis de \textit{in-play} (en vivo). A nivel pre-match, reinan variaciones de sistemas de Elo/TrueSkill.
    
    \item \textbf{¿Qué métrica domina?} La literatura frecuentista está obsesionada con el \textbf{Accuracy} (Precisión) y el AUC-ROC. Es preocupante que, en un campo (apuestas y estrategia deportiva) donde el grado de certidumbre es vital, las métricas estrictamente orientadas a la evaluación probabilística como el Brier Score o las curvas de calibración aparecen como secundarias.
    
    \item \textbf{Sesgos del campo / ¿Qué nadie cuestiona?} La "ilusión de contingencia del roster". Prácticamente todos los modelos analizados entrenan con datos agregados por el \textit{"-tag-"} o nombre orgánico del club (ej. FaZe Clan, Natus Vincere) asumiendo su estacionariedad temporal, obviando que en CS:GO los cinco jugadores son los que dictan el "TrueSkill", no la marca.
    
    \item \textbf{Vacíos argumentados:} Hay una oportunidad clara para un modelo que logre un equilibrio: la \textbf{interpretabilidad} de una regresión logística, combinada metodológicamente con variables medioambientales (Mapa elegido, descanso) que TrueSkill ignora, y evaluada con el rigor de la \textbf{incertidumbre epistémica} que los modelos de Deep Learning son incapaces de ofrecer.
    
    \item \textbf{Implicación fundamental para la Pregunta de Investigación:} Este mini estado del arte justifica sólidamente la decisión del proyecto. La estimación bayesiana para probabilidades pre-match soluciona directamente el vacío de \textit{overconfidence} (exceso de confianza de los modelos frecuentistas) y ofrece probabilidades reales y creíbles (HDi), un diferencial gigantesco frente a enfoques estáticos.
\end{enumerate}

\end{document}
